\usepackage{graphicx}
\graphicspath{ {./figures/} } % Figure는 figures 폴더에 저장

\usepackage{csquotes}

\usepackage{wrapfig}
\usepackage{amsmath}
\usepackage{subcaption}
\usepackage{hyperref}
\usepackage{tikz}
\usepackage{tikzscale}
\usepackage{tabularx}
\usepackage{indentfirst}

% ---- 이 보고서에서 추가한 설정 ----
\usepackage{booktabs,array,float,enumitem}
\usepackage{environ}
\usepackage[most]{tcolorbox}
\setlist{nosep,leftmargin=16pt,itemsep=1pt}
\newcolumntype{L}[1]{>{\raggedright\arraybackslash}p{#1}}

% 캡션: 청소년 과학창의연구 양식 ([Fig. 1], Table 1.)
\DeclareCaptionLabelFormat{figbr}{[Fig.~#2]}
\captionsetup{font=small,labelfont=bf,justification=justified}
\captionsetup[figure]{labelformat=figbr,labelsep=space}
\captionsetup[table]{name=Table,labelsep=period}
\renewcommand{\figurename}{Fig.}

% =====================================================================
% 중복 식 블록
%   같은 내용을 두 가지 방법으로 쓴 곳이다.
%   dupA = 더 기본적이고 직관적인 유도 (먼저)
%   dupB = 이론 문서(2026.10.31)의 원래 식, 그리고 dupA와 같다는 확인
%   아래 \dupchoose 값을 바꾸면 고를 수 있다.
%     both : 둘 다 보이고, 주황 상자로 표시 (지금)
%     A    : 가만 남김 (상자 없이 본문처럼)
%     B    : 나만 남김 (상자 없이 본문처럼)
% =====================================================================
\newcommand{\dupchoose}[2]{\csdef{dupsel-#1}{#2}}
\dupchoose{D1}{both}   % 속도 연결: 체인 속도 / 크랭크 각속도
\dupchoose{D2}{both}   % 모멘트팔: 수직거리 / 극좌표 공식
\dupchoose{D3}{both}   % 페달 토크: 힘과 팔 / 일의 관계
\dupchoose{D4}{both}   % 한 바퀴의 일: 에너지 / 등속 토크식 적분
\dupchoose{D5}{both}   % 최적 모양: 같은 이득 조건 / 코시-슈바르츠
\dupchoose{D6}{both}   % 체인 장력: 앞축 식 (3) / 다리 관성 포함 식

\definecolor{dupc}{HTML}{C2650F}
\tcbset{dupbox/.style={breakable,enhanced,colback=white,colframe=dupc,boxrule=0.7pt,arc=0pt,
  left=5pt,right=5pt,top=2pt,bottom=2pt,fonttitle=\small\bfseries,coltitle=white,
  before skip=6pt,after skip=6pt}}
\newcommand{\dupSel}[1]{\ifcsundef{dupsel-#1}{both}{\csuse{dupsel-#1}}}
\NewEnviron{dupA}[2]{%
  \ifcsundef{dupsel-#1}{\def\dupS{both}}{\letcs\dupS{dupsel-#1}}%
  \ifdefstring{\dupS}{both}{%
    \begin{tcolorbox}[dupbox,colbacktitle=dupc,title={중복 #1-가 \quad 기본·직관적 유도: #2}]\BODY\end{tcolorbox}}{%
  \ifdefstring{\dupS}{A}{\BODY\par}{}}}
\NewEnviron{dupB}[2]{%
  \ifcsundef{dupsel-#1}{\def\dupS{both}}{\letcs\dupS{dupsel-#1}}%
  \ifdefstring{\dupS}{both}{%
    \begin{tcolorbox}[dupbox,colbacktitle=dupc!65!black,title={중복 #1-나 \quad 이론 문서의 식과 같음: #2}]\BODY\end{tcolorbox}}{%
  \ifdefstring{\dupS}{B}{\BODY\par}{}}}

% 참고문헌 글자 크기
\renewcommand*{\bibfont}{\small}
\setlength{\bibitemsep}{0pt}

% 그림이 빈칸 없이 들어가도록
\renewcommand{\topfraction}{0.9}
\renewcommand{\bottomfraction}{0.8}
\renewcommand{\textfraction}{0.07}
\renewcommand{\floatpagefraction}{0.8}
\setcounter{totalnumber}{4}
